\section{Post-hoc support-conditioned predictor diagnostics}
\label{app:support_conditioned_predictors}

This appendix records the post-hoc support-conditioned predictor experiments. These experiments are diagnostic: they ask whether learned support families can select fitted transition maps after the Koopman autoencoder has already been trained. They are not used for the main paper's basin-identifiability claim or for the main forecasting table.

\paragraph{Support-conditioned rollouts.}
Support-conditioned prediction is a second-stage procedure applied after the Koopman autoencoder has been trained with a single global transition \(K\). In the second stage, the encoder \(\Enc\), decoder \(\Dec\), and original \(K\) are frozen. A separate centered transition \(K_c\) is optimized for each retained \(F_{\rm top8}\) support family \(c\), while the original \(K\) remains the fallback for unretained or unassigned routes. The comparison point is the same frozen checkpoint's global rollout,
\begin{equation}
  \hat z_{t+1}^{\rm global}=K\hat z_t,\qquad
  \hat x_{t+1}^{\rm global}=\Dec(\hat z_{t+1}^{\rm global}).
  \label{eq:global_rollout}
\end{equation}

Routes are constructed without basin labels. On disjoint fitting trajectories, we encode states with the frozen encoder, compute exact \(S_{\rm top8}\) supports, and merge them into support families by the Jaccard rule in \cref{sec:method_supports} with threshold \(0.40\). Families with fewer than \(50\) current-state fitting transitions are not given trainable maps. For each retained family \(c\in\mathcal C_{\rm fit}\), the fixed center \(\bar z_c\) is the mean current latent assigned to that family, and the map is initialized from the checkpoint transition, \(K_c\leftarrow K\). The route-local transition is
\begin{equation}
  T_c(z)=\bar z_c+K_c(z-\bar z_c).
  \label{eq:family_local_map}
\end{equation}

The reported local maps are trained jointly as one bank of matrices by gradient descent on decoded rollout MSE. For a second-stage training window \((x_{b,0},\ldots,x_{b,H})\), let \(\hat z_{b,0}=\Enc(x_{b,0})\), where \(H\) is the checkpoint's short training horizon. At rollout offset \(\ell\), the current predicted latent is detached only for the discrete route assignment:
\[
  c_{b,\ell}=F_{\rm top8}(\hat z_{b,\ell}).
\]
If this assigned family is retained, the next latent uses \(T_{c_{b,\ell}}\); otherwise it falls back to the frozen global transition:
\begin{equation}
  \tilde z_{b,\ell+1}=
  \begin{cases}
    T_{c_{b,\ell}}(\hat z_{b,\ell}), & c_{b,\ell}\in\mathcal C_{\rm fit},\\
    K\hat z_{b,\ell}, & c_{b,\ell}\notin\mathcal C_{\rm fit}.
  \end{cases}
  \label{eq:family_local_rollout}
\end{equation}
The decoded prediction is \(\hat x_{b,\ell+1}=\Dec(\tilde z_{b,\ell+1})\). Every \(m\) predicted steps, with \(m=5\) in the selected runs, the next latent state is refreshed by the frozen encoder,
\[
  \hat z_{b,\ell+1}=\Enc(\hat x_{b,\ell+1}),
\]
and otherwise \(\hat z_{b,\ell+1}=\tilde z_{b,\ell+1}\). The active setting recomputes \(c_{b,\ell}\) before every one-step map application, so route switching can occur at any predicted step, not only at refresh boundaries.

The second-stage objective is
\begin{equation}
  \min_{\{K_c:\,c\in\mathcal C_{\rm fit}\}}
  \frac{1}{BH}\sum_{b=1}^{B}\sum_{\ell=1}^{H}
  \norm{\hat x_{b,\ell}-x_{b,\ell}}_2^2 .
  \label{eq:stage2_local_loss}
\end{equation}
There is no basin-label loss, support-classification loss, latent consistency loss, or regularizer in this stage. Minibatches are route-balanced: a training step samples support-family buckets approximately uniformly and then samples windows inside the selected buckets. Thus frequent support families do not monopolize the updates, and a particular \(K_c\) receives gradients only on examples and rollout steps where it is selected.

At evaluation time no parameters are updated. Starting from a held-out initial state, the same recurrence in \cref{eq:family_local_rollout} is run for the forecast horizon using only predicted states. Exact supports seen in the codebook map directly to their family; unseen exact supports are assigned to the nearest family prototype only when the Jaccard similarity reaches the same threshold, and otherwise use the global fallback. Exact \(S_{\rm top8}\)-local routing uses the same rollout structure with exact supports as routes and is retained as an ablation in \cref{app:routing_robustness}.

\begin{table}[tbp]
\centering
\caption{Appendix support-conditioned predictor diagnostic. Entries are label-free \(F_{\rm top8}\)-local routed/global MSE ratios and system-win counts on all states; lower ratios and higher win counts are better. Superscripts denote system-level significance tests: \(\ast\), \(p_{\rm Holm}<0.05\); \(\ast\ast\), \(p_{\rm Holm}<0.01\).}
\label{tab:self_routing}
\setlength{\tabcolsep}{2pt}
\resizebox{0.78\textwidth}{!}{\begin{tabular}{@{}l cc cc@{}}
\toprule
& \multicolumn{2}{c}{$H100$} & \multicolumn{2}{c}{$H1000$} \\
\cmidrule(lr){2-3}\cmidrule(l){4-5}
Model & \(F_{\rm top8}\) ratio & system wins & \(F_{\rm top8}\) ratio & system wins \\
\midrule
LISTA & ${0.239}^{\ast}$ & $13/15$ & ${\mathbf{1.41\times10^{-11}}}^{\ast\ast}$ & $15/15$ \\
LISTA-SB & ${0.021}^{\ast}$ & $13/15$ & ${2.42\times10^{-7}}^{\ast}$ & $14/15$ \\
LISTA-BD & ${\mathbf{0.011}}^{\ast}$ & $13/15$ & ${6.96\times10^{-5}}^{\ast}$ & $15/15$ \\
Sparse MLP, BD & ${0.032}^{\ast\ast}$ & $12/15$ & ${6.53\times10^{-8}}$ & $15/15$ \\
Sparse MLP & ${0.038}^{\ast\ast}$ & $13/15$ & ${5.19\times10^{-4}}$ & $15/15$ \\
Dense MLP & ${154}$ & $1/15$ & ${85.7}$ & $1/15$ \\
\bottomrule
\end{tabular}}
\end{table}
